\documentclass[11pt]{article}

\usepackage[margin=1.7cm]{geometry}
\usepackage{amsmath}
\usepackage{amssymb}
\usepackage{enumitem}
\usepackage{titlesec}
\usepackage{xcolor}
\usepackage{fancyhdr}
\usepackage{multicol}
\usepackage{array}

\titleformat{\section}{\large\bfseries\color{blue!50!black}}{\thesection}{0.6em}{}
\titlespacing{\section}{0pt}{1.1em}{0.4em}
\titleformat{\subsection}{\normalsize\bfseries}{\thesubsection}{0.6em}{}
\titlespacing{\subsection}{0pt}{0.8em}{0.3em}

\setlength{\headheight}{14pt}

\pagestyle{fancy}
\fancyhf{}
\fancyhead[L]{Ratios in Algebra}
\fancyfoot[C]{\thepage}
\renewcommand{\headrulewidth}{0.4pt}

\fancypagestyle{firstpage}{%
  \fancyhf{}
  \fancyhead[L]{Ratios in Algebra}
  \fancyhead[R]{Name: \makebox[4cm]{\hrulefill}}
  \fancyfoot[C]{\thepage}
  \renewcommand{\headrulewidth}{0.4pt}
}

\newcommand{\ans}{\hrulefill}
\newcommand{\anshere}{\par\vspace{0.2em}\noindent\ans\par\vspace{0.2em}}
\setlist[enumerate,1]{leftmargin=1.4em, itemsep=0.4em, topsep=0.2em}
\setlength{\columnsep}{1.1cm}
\setlength{\columnseprule}{0.4pt}
\raggedcolumns

\begin{document}
\thispagestyle{firstpage}

\begin{center}
  {\Large\bfseries Ratios in Algebra}
\end{center}
\vspace{0.2em}

\section*{Fluency}

\begingroup\small
\begin{multicols}{2}
\subsection*{Simplify}
\noindent\textbf{Worked example:} Simplify $12xy:18y$. Divide both sides by the HCF,
$6y$: $12xy:18y = 2x:3$.
\vspace{0.3em}
\begin{enumerate}[label=(\alph*)]
  \item $6x:8x$
  \anshere
  \item $10x:15x$
  \anshere
  \item $4x:12x$
  \anshere
  \item $9x:6x$
  \anshere
  \item $5xy:10y$
  \anshere
  \item $8x^2:4x$
  \anshere
  \item $3x:9x^2$
  \anshere
  \item $12ab:18b$
  \anshere
\end{enumerate}

\columnbreak

\subsection*{Find the unknown}
\noindent\textbf{Worked example:} If $x:y=2:5$ and $x=8$, find $y$. Since
$8\div 2=4$, multiply the other side by $4$ too: $y=5\times 4=20$.
\vspace{0.3em}
\begin{enumerate}[label=(\alph*)]
  \item $x:y=3:4$. If $x=15$, find $y$.
  \anshere
  \item $a:b=5:2$. If $b=8$, find $a$.
  \anshere
  \item $p:q=2:7$. If $q=21$, find $p$.
  \anshere
  \item $m:n=9:4$. If $m=27$, find $n$.
  \anshere
  \item $x:y=7:3$. If $y=18$, find $x$.
  \anshere
  \item $a:b=11:5$. If $a=55$, find $b$.
  \anshere
\end{enumerate}
\end{multicols}
\endgroup

\newpage

\section*{Problem Solving}

\begingroup\small
\begin{multicols}{2}
\subsection*{Solve}
\noindent\textbf{Worked example:} Solve $(x+2):(x-1)=5:3$. Cross-multiply, then
solve the equation:
\begin{align*}
  3(x+2) &= 5(x-1) \\
  3x+6 &= 5x-5 \\
  11 &= 2x \\
  x &= 5.5.
\end{align*}
\vspace{0.2em}
\begin{enumerate}[label=(\alph*)]
  \item $(x+3):(x-1)=3:2$
  \anshere
  \item $(3x+1):(x+7)=2:1$
  \anshere
\end{enumerate}

\vspace{0.4em}
\subsection*{Sharing in a ratio}
\noindent\textbf{Worked example:} Share \$80 in the ratio $3:5$. The parts are
$3k$ and $5k$, so $3k+5k=80 \Rightarrow k=10$. The shares are $3\times10=\$30$ and
$5\times10=\$50$.
\vspace{0.2em}
\begin{enumerate}[label=(\alph*)]
  \item A $40$\,cm line segment is divided in the ratio $3:5$. Find the length of
    each part.
  \anshere
  \item Three siblings share \$780 in the ratio $2:3:8$. How much does each
    receive?
  \anshere
\end{enumerate}

\columnbreak

\subsection*{Word problems}
\noindent\textbf{Worked example:} $a:b=5:3$ and $a-b=14$. Find $a$ and $b$. Write
$a=5k$ and $b=3k$, so $5k-3k=14 \Rightarrow k=7$. Then $a=35$ and $b=21$.
\vspace{0.2em}
\begin{enumerate}[label=(\alph*)]
  \item The ratio of cats to dogs at a shelter is $3:4$. There are $12$ more dogs
    than cats. How many of each are there?
  \anshere
  \item Henry and Jack share some money in the ratio $3:5$. Henry receives \$45.
    How much does Jack receive, and what is the total amount shared?
  \anshere
\end{enumerate}
\end{multicols}
\endgroup

\newpage

\section*{Reasoning}
\noindent\textit{No algebra here --- just ratio reasoning. Show your working.}

\begin{enumerate}[label=(\alph*)]
  \item Explain why the ratios $8:12$ and $2:3$ represent the same relationship.
  \anshere
  \item Cordial A is mixed $1$ part cordial to $4$ parts water. Cordial B is mixed
    $2$ parts cordial to $9$ parts water. Which is the stronger-tasting cordial?
    Justify your answer using equivalent fractions or ratios.
  \anshere
  \item A fruit drink is made from orange and mango juice in the ratio $5:3$.
    Another is made in the ratio $7:4$. Which drink has the higher proportion of
    orange juice? Justify your answer.
  \anshere
  \item Paint colour ``Sunrise'' is mixed from red and yellow paint in the ratio
    $3:7$. You want to make a larger batch using $18$\,L of red paint. How much
    yellow paint do you need?
  \anshere
  \item A map has a scale of $1:25\,000$. Explain what this ratio means, then
    calculate the real-world distance, in km, represented by $6$\,cm on the map.
  \anshere
\end{enumerate}

\section*{Enrichment}

\subsection*{Combining Ratios}
\noindent\textit{To combine $a:b$ and $b:c$ into one ratio $a:b:c$, scale each
ratio so that the value of $b$ matches in both.}

\vspace{0.3em}
\noindent\textbf{Worked example:} $a:b=2:3$ and $b:c=4:5$. The lowest common
multiple of $3$ and $4$ is $12$, so scale the first ratio by $4$ and the second by
$3$: $a:b=8:12$ and $b:c=12:15$. Combining these, $a:b:c=8:12:15$.

\vspace{0.3em}
\begin{enumerate}[label=(\alph*)]
  \item $a:b=3:4$ and $b:c=2:5$. Find $a:b:c$.
  \anshere
  \item $x:y:z=2:5:9$. If $z-x=42$, find $x$, $y$ and $z$.
  \anshere
\end{enumerate}

\subsection*{A Harder Ratio Problem}
\noindent\textbf{Worked example:} Three years ago, the ratio of Mia's age to her
brother's age was $5:2$. In two years' time, the ratio will be $2:1$. Find their
current ages. Let Mia's current age be $M$ and her brother's be $B$.
\[\text{3 years ago: } \frac{M-3}{B-3}=\frac{5}{2} \;\Rightarrow\; 2(M-3)=5(B-3)
  \;\Rightarrow\; M = \frac{5B-9}{2}.\]
\[\text{In 2 years: } \frac{M+2}{B+2}=\frac{2}{1} \;\Rightarrow\; M+2=2(B+2)
  \;\Rightarrow\; M=2B+2.\]
Setting these equal: $\dfrac{5B-9}{2}=2B+2 \Rightarrow 5B-9=4B+4 \Rightarrow B=13$,
so $M=2(13)+2=28$.

\vspace{0.3em}
\begin{enumerate}[label=(\alph*)]
  \item Six years ago, the ratio of Ana's age to Ben's age was $3:1$. In four
    years' time, the ratio will be $2:1$. Find their current ages.
  \anshere
\end{enumerate}

\end{document}
